\section{Progress Tracking}\label{sec:id-progress}

Progress is recorded for each objective in Table~\ref{tab:objectives}. The record is kept on the student's own computer, as described in the Data and Persistence section of the SE document. It stores objective counts and recent answers for mastery, as well as operation-category counters from the tree activity.

\subsection{What Is Recorded}
For each of the five objectives shown on the progress screen, the app records attempts, first-answer correctness, hints used, and the recent answers needed to calculate mastery status. The record does not store a last-update time. These counts allow the app to distinguish a student who answers correctly without help from one who answers correctly only after a hint.

In the tree activity, the first answer to each step is recorded under an operation-category key. Predict-card answers use \texttt{predict-insert}, \texttt{predict-search}, \texttt{predict-delete-copy} or \texttt{predict-delete-merge}. Answers on Build-card traces use \texttt{trace-insert} for insertion and \texttt{trace-search} for search and deletion. The generic tracker accepts these keys and persists their counts, but the progress screen displays only the five identifiers in Table~\ref{tab:objectives}; operation-category counts are not mapped to those objectives or shown as objective mastery. Undo and Redo do not record answers.

\subsection{Mastery Standard}
An objective is shown as \emph{practised} once the student has answered at least five items on it. It is shown as \emph{secure} when at least 80\% of the student's last five first answers were correct and none of them used a hint. The standard follows the mastery learning approach \cite{Bloom1968}, in which a student is held to a defined standard rather than ranked against other students.

\subsection{What the Student Sees}
The progress screen shows one bar per objective with its identifier, title and description. Status can be untried, started, practised or secure. It also shows attempts, first-time correct answers and hints used, and states the mastery standard in words. The screen recommends the first lesson whose objective is not secure, but this recommendation does not lock access to other lessons.

The progress screen supports self-regulation, a recurring issue in the module and the focus of UDL Guideline 9 (self-regulation)~\cite{CAST2018}: learners can see what they have practised and what remains, and plan their next step. In Kolb's terms~\cite{Kolb1984} it is reflective observation. The screen recommends a lesson for further practice but does not route the learner automatically; students can choose whether to review.
 
The interface labels this count ``Hints used''; it does not rename it ``support used''.

\subsection{Limits of the Record}
The record shows what a student has done in the app. It does not show whether they learned the material in the lecture, or whether the record reflects a single good session. The evaluation in Section~\ref{sec:id-evaluation} is designed to test whether the record is a useful guide.